\section{The reduced fold jet and transverse return}
\label{sec:fold-jet}

We now identify the coefficient that distinguishes the two delay
organizations.  All formulas in this section are coefficients of the actual
history graph from \cref{thm:history-graph}; the symbolic calculation is a
check, not a substitute for the graph theorem.

\subsection{Exact chart equations}

For a scalar history \(\phi\), write \(\phi_j=\phi(s-\theta_j)\) and set
\begin{equation}
\label{eq:delay-operators}
\begin{aligned}
 \cD_c\phi&=\phi-\frac13\phi_0-\frac23\phi_1,\\
 \cD_{cz}\phi&=-\frac16\phi+\frac1{12}\phi_0+\frac1{12}\phi_1,
 \qquad \cD_z\phi=-\cD_{cz}\phi,\\
 \Delta\phi&=\phi_0-\phi_1.
\end{aligned}
\end{equation}
Substitution of \cref{eq:blowup-scaling} into
\cref{eq:physical-model}, followed by projection onto the biorthogonal
bases, gives the exact polynomial RFDE
\begin{equation}
\label{eq:exact-chart}
\begin{aligned}
X'={}&Y-\alpha X^2
 +\del\left[K\cD_cX-2\alpha XZ-\frac{20}{9}\alpha^2X^3\right]\\
&+\del^2\left[K\cD_{cz}Z-\alpha Z^2+4\alpha^2X^2Z\right]
-\del^3\frac{20}{3}\alpha^2XZ^2
+\del^4\frac43\alpha^2Z^3,\\
Y'={}&-X+\del\nu,\\
\del Z'={}&-2Z-\alpha X^2
+\del\left[-2\alpha XZ+\frac43\alpha^2X^3-K\eta\Delta X\right]\\
&+\del^2\left[-W-\alpha Z^2-\frac{20}{3}\alpha^2X^2Z
              +K(\cD_zZ-\eta\Delta Z)\right]\\
&+\del^3 4\alpha^2XZ^2-\del^4\frac{20}{9}\alpha^2Z^3,\\
\del W'={}&Z-d_wW.
\end{aligned}
\end{equation}
In particular, setting \(\del=0\) gives
\cref{eq:singular-stable-graph,eq:singular-critical-system}.

For later use, let \(\phi_0^s\) denote the complete flow of the prepared
singular field and put
\begin{equation}
\label{eq:singular-backtracks}
 X_{-j}(u)=\pi_X\phi_0^{-\theta_j}(u),
 \qquad \Delta_0X=X_{-0}-X_{-1}.
\end{equation}

\begin{proposition}[First graph coefficients]
\label{prop:transverse-return}
On every uncut portion of the logarithmic core and its depth-two flow hull,
the first reduced coefficient in \cref{eq:growing-graph-expansion} is
\begin{equation}
\label{eq:Q1}
 Q_1(X,Y)=\vect{
 K\left(X-\frac13X_{-0}-\frac23X_{-1}\right)
 -\frac{11}{9}\alpha^2X^3\\ \nu}.
\end{equation}
The leading stable forcing after the shift \cref{eq:stable-shift} is
\begin{equation}
\label{eq:G0}
 G_0(X,Y)=\vect{
 \alpha XY+\frac43\alpha^2X^3-K\eta\Delta_0X\\
 \dfrac{\alpha}{d_w}X(Y-\alpha X^2)}.
\end{equation}
Consequently
\begin{equation}
\label{eq:H1-eta}
 \partial_\eta H_{1,\widehat Z}=-\frac K2\Delta_0X,
 \qquad
 \partial_\eta H_{1,\widehat W}=-\frac{K}{2d_w}\Delta_0X.
\end{equation}
Along the singular canard,
\begin{equation}
\label{eq:delayed-canard-difference}
 \Delta_0X_0=\frac{\theta_0-\theta_1}{2\alpha},
\end{equation}
and the first nonzero return to the critical equation is
\begin{equation}
\label{eq:transverse-return-coefficient}
 \partial_\eta Q_{2,X}(\gamma_0(s),\nu,0)
 =-\frac{K(\theta_0-\theta_1)}{4\alpha}s.
\end{equation}
Moreover,
\begin{equation}
\label{eq:transverse-return-y-zero}
 \partial_\eta Q_{2,Y}(\gamma_0(s),\nu,0)=0,
\end{equation}
because the exact critical recovery equation is
\(Y'=-X+\del\nu\) and contains no delayed or transverse term.
In particular this coefficient is independent of \(d_w\).
\end{proposition}

\begin{proof}
Insert \cref{eq:stable-shift} into \cref{eq:exact-chart}.  The stable
linear part becomes the matrix \(A\) in \cref{eq:q0-A}; every remaining
term is divisible by \(\del\).  Differentiating the fixed-point equations
\cref{eq:graph-transform} once at \(\del=0\) gives
\begin{equation}
 Q_1=F(\cE_{q_0,0};0,\nu,\eta),
 \qquad H_1=-A^{-1}G(\cE_{q_0,0};0,\nu,\eta).
\end{equation}
Direct collection of the critical terms gives \cref{eq:Q1}; collection of
the shifted transverse terms gives \cref{eq:G0}.  Since
\[
 -A^{-1}=\mat{1/2&0\\1/(2d_w)&1/d_w},
\]
differentiating with respect to \(\eta\) yields
\cref{eq:H1-eta}.

The singular orbit has \(X_0(s)=-s/(2\alpha)\), hence
\[
 X_0(s-\theta_0)-X_0(s-\theta_1)
 =\frac{\theta_0-\theta_1}{2\alpha},
\]
which proves \cref{eq:delayed-canard-difference}.  The coefficient \(Q_1\)
is independent of \(\eta\).  At the next order the only first
\(\eta\)-derivative returning to the critical equation is the product of
the critical variable with the transverse graph displacement:
\[
 \partial_\eta Q_{2,X}
 =-2\alpha X\,\partial_\eta H_{1,\widehat Z}
 =\alpha KX\Delta_0X.
\]
Evaluating on \(\gamma_0\) gives
\cref{eq:transverse-return-coefficient}.  The \(\widehat W\) component
does not enter this order, which explains the absence of \(d_w\).
\end{proof}

The mechanism can now be read directly from the equations:
\[
 \text{fixed critical delay measure}
 \longrightarrow
 \text{transverse delayed difference}
 \longrightarrow
 \partial_\eta H_{1,\widehat Z}
 \longrightarrow
 \partial_\eta Q_{2,X}.
\]
The remaining task is to show that the last quantity shifts an actual
complete-history connection root.  That requires a one-sided selection that
removes both the time-translation mode and the Gaussian normal mode.
